\section{Replicação do \modus\ (Etapa 1) e Vazamento de Rótulo}
\label{sec:replication}

Não há código público dos autores. Reimplementamos o \modus\ a partir do texto, em Python, com o
HST e o ADWIN do CapyMOA~\cite{gomes2025capymoa}, que encapsula o MOA~\cite{bifet2010moa}. Os
hiperparâmetros são os da Seção~4.1 do artigo, por base. Nas três bases chegamos exatamente ao
número de instâncias reportado. No ASEADOS são 312.262 fluxos (classes Benign e DDoS). No CTU-13
são 6.336.394 pacotes IPv4 do \emph{pcap} truncado do cenário 11. No Edge-IIoTset são 30.919.961
quadros: a única combinação que fecha esse total usa os \emph{pcaps} normais duas vezes, o
\emph{Port Scanning} como benigno e os quatro DDoS.

\subsection{Resultados da Etapa 1}
A Tabela~\ref{tab:repl} resume a replicação. No ASEADOS e no Edge-IIoT os números do artigo se
reproduzem; no Edge-IIoT, os verdadeiros positivos e falsos negativos são idênticos pacote a pacote.
No CTU-13 a replicação não fecha: os rótulos do artigo são internamente inconsistentes. A regra
oficial do CTU-13 (IPs infectados) dá 32,04\% de ataques, exatamente o valor que a Tabela~1 do
artigo chama de ``Benigno'', e a matriz de confusão implica 2,89 milhões de ataques, total que
nenhuma regra testada reproduz. A análise da Etapa~1 mostrou ainda que:
(i) a \emph{flag} do ADWIN, contribuição central do artigo, não altera nenhuma predição no CTU-13 e
altera um pacote em 19,6 milhões no Edge-IIoT;
(ii) sem normalização, o HST separa por protocolo e não por volume: os cortes caem em $[0,1]$ e
portas e tamanhos ficam sempre do mesmo lado. O recall de 73,16\% no Edge-IIoT é exatamente a
fração de DDoS que não é TCP;
(iii) o limiar passa quase todo o tempo preso no piso $\tau_{\min}$, ajustado à mão por base;
(iv) no ASEADOS, a ordem ``temporal'' que reproduz o artigo vem de um erro de leitura de datas
(dd/mm lido como mm/dd). Na ordem cronológica correta o DDoS ocorre durante o aquecimento e o F1 é 0\%.

\begin{table}[t]
\centering
\caption{Replicação do \modus\ (F1 \%, convenção do artigo). ``Com rótulo'': o HST do MOA recebe o
rótulo verdadeiro como atributo (Etapa~1). ``Sem rótulo'': fluxo corrigido. $^{*}$Medido no conjunto
sem a duplicação do tráfego normal (19,6 milhões de quadros), no qual a rodada com rótulo também dá 84,42.}
\label{tab:repl}
\begin{tabular}{lccc}
\toprule
Base & Artigo & Com rótulo & Sem rótulo \\
\midrule
ASEADOS-SDN-IoT & 99,62 & 99,73 & \textbf{0,00} \\
CTU-13 (C11)    & 83,79 & 64,55 & 64,55 \\
Edge-IIoTset    & 81,24 & 84,33 & 84,42$^{*}$ \\
\bottomrule
\end{tabular}
\end{table}

\subsection{O HSTrees do MOA usa o rótulo como atributo}
\label{sec:leak}
Ao portar o HST para \texttt{numpy} (Seção~\ref{sec:hstnp}), notamos que o MOA dimensiona o espaço de
trabalho com \texttt{instance.numAttributes()}, que \emph{inclui o atributo de classe}. Quando o fluxo
carrega o rótulo verdadeiro, como no uso padrão do CapyMOA, em que \texttt{NumpyStream} recebe
\texttt{X} e \texttt{y}, cerca de $1/(d+1)$ dos cortes de cada árvore separam ataque de benigno pelo
próprio rótulo. Isso também vale para a pontuação, porque \texttt{score\_instance} recebe a mesma
instância. Verificamos o efeito diretamente: com o mesmo vetor de atributos e só o rótulo trocado,
o score do MOA muda até 0,10.

O efeito sobre a replicação depende de os atributos já separarem as classes. No CTU-13 eles
separam: o \emph{flood} é ICMP. As predições com e sem rótulo são idênticas nas 6,3 milhões de
instâncias. No ASEADOS não separam: todo fluxo DDoS é o vetor $[0,0,0,0]$, idêntico a 42.740 fluxos
benignos, e o que dava a esses fluxos o score constante de 0,806, acima do piso de 0,75, era a
dimensão do rótulo. Sem ela, o recall cai de 100\% para 0\% e a AUC-ROC de 1,00 para 0,11. O
detector passa a considerar o DDoS \emph{mais normal} que o tráfego benigno: o ataque passa no
início, treina o modelo e vira o perfil de referência.

Não sabemos se a implementação dos autores tem o mesmo problema, pois o código não é público. O
resultado mostra, porém, que os 99,62\% no
ASEADOS são reproduzíveis \emph{exatamente} com o vazamento e não são reproduzíveis sem ele. Daqui
em diante, todas as rodadas usam o fluxo corrigido: o MOA recebe um rótulo constante, e o motor
\texttt{numpy} não tem a dimensão do rótulo.
